\documentclass[12pt]{article}
\usepackage[utf8]{inputenc}
\usepackage[T1]{fontenc}
\usepackage[margin=1in]{geometry}
\usepackage{setspace}
\usepackage{xcolor}
\usepackage[colorlinks=true, urlcolor=blue!60!black, linkcolor=blue!60!black]{hyperref}
\setlength{\parindent}{2em}
\setlength{\parskip}{1.25em}
\renewcommand{\baselinestretch}{1.0}
\title{\textit{Trump v. Anderson}, 601 U.S. 100 (2024)}
\author{Case Brief}
\date{}
\begin{document}
\maketitle

\section{Case Metadata}

\begin{itemize}
  \item \textbf{Case Name:} \emph{Trump v. Anderson}
  \item \textbf{Citation:} \href{https://www.supremecourt.gov/opinions/23pdf/601us1r06_m6hn.pdf}{601 U.S. 100 (2024)} (Docket No.\ 23--719)
  \item \textbf{Court:} Supreme Court of the United States (on writ of certiorari to the Supreme Court of Colorado)
  \item \textbf{Dates:} Argued February 8, 2024; Decided March 4, 2024
  \item \textbf{Vote:} 9--0 judgment reversing the Colorado Supreme Court; 5--4 division on whether congressional legislation under Section~5 is the exclusive federal mechanism for enforcing Section~3
  \item \textbf{Opinion Authorship:} Per Curiam (joined in full by Roberts, C.J., and Thomas, Alito, Gorsuch, and Kavanaugh, JJ.); Barrett, J., concurring in part (joining Parts I and II--B) and concurring in the judgment; Sotomayor, Kagan, and Jackson, JJ., jointly concurring in the judgment
  \item \textbf{Briefing Book Placement:} Chapter 4: Federalism and State Sovereignty
  \item \textbf{Primary Texts:} \href{https://www.supremecourt.gov/opinions/23pdf/23-719_19m2.pdf}{SupremeCourt.gov Slip Opinion} $\mid$ \href{https://www.law.cornell.edu/supremecourt/text/23-719}{Cornell LII} $\mid$ \href{https://supreme.justia.com/cases/federal/us/601/23-719/}{Justia}
\end{itemize}

\section{Facts and Procedural History}

In September 2023, six months before Colorado's March 5, 2024 presidential primary, four Republican and two unaffiliated Colorado voters---led by Norma Anderson---petitioned Colorado state court under the Colorado Election Code (Colo.\ Rev.\ Stat.\ \S\S 1--1--113(1), 1--4--1101(1), 1--4--1201, 1--4--1203(2)(a), 1--4--1204) against former President Donald J.\ Trump and Colorado Secretary of State Jena Griswold (\href{https://www.law.cornell.edu/supremecourt/text/23-719}{Cornell LII}). Respondents alleged that after losing the 2020 presidential election, Trump intentionally organized and incited the crowd that breached the U.S.\ Capitol on January 6, 2021, thereby disqualifying himself from again holding the Presidency under Section~3 of the Fourteenth Amendment.

Following a five-day trial, Denver District Court Judge Sarah B.\ Wallace issued an order on November 17, 2023, finding that Trump had ``engaged in insurrection'' within the meaning of Section~3, yet denying the petition because the Presidency is not an ``office \ldots\ under the United States'' and the President is not an ``officer of the United States'' covered by Section~3 (\href{https://law.justia.com/cases/colorado/supreme-court/2023/23sa300.html}{\emph{Anderson v. Griswold}, 2023 CO 63}).

On December 19, 2023, the Colorado Supreme Court reversed in part and affirmed in part by a 4--3 vote in \href{https://law.justia.com/cases/colorado/supreme-court/2023/23sa300.html}{\emph{Anderson v. Griswold}, 2023 CO 63, 543 P.3d 283}. The majority held that Section~3 encompasses the Presidency and the President, is self-executing without congressional legislation, is justiciable under the Colorado Election Code, and barred Trump from the 2024 Republican primary ballot; it ordered Secretary Griswold not to list Trump's name or count write-in votes for him. Chief Justice Boatright and Justices Samour and Berkenkotter each dissented separately, raising statutory-scope, due-process, and enforcement objections. The Colorado Supreme Court stayed its ruling pending U.S.\ Supreme Court review. On January 5, 2024, the U.S.\ Supreme Court granted certiorari on an expedited schedule, heard argument on February 8, and issued its decision on March 4, 2024---one day before Super Tuesday---ordering the mandate to issue forthwith (\href{https://www.supremecourt.gov/opinions/23pdf/601us1r06_m6hn.pdf}{601 U.S. at 104--106, 117}).

\section{Issue(s) Presented}

\begin{enumerate}
  \item Did the Colorado Supreme Court err in ordering former President Donald J.\ Trump excluded from Colorado's 2024 Republican presidential primary ballot under Section~3 of the Fourteenth Amendment?
  \item More broadly, do individual States possess constitutional authority to enforce Section~3 against federal officeholders and candidates, or does enforcement authority over federal offices rest exclusively with Congress under Section~5 of the Fourteenth Amendment?
\end{enumerate}

\section{Holding}

\textbf{Reversed (9--0 judgment; 5--4 on Section~5 exclusivity).} All nine Justices held that States lack constitutional power to enforce Section~3 of the Fourteenth Amendment against federal officeholders and candidates---especially presidential candidates---although States retain sovereign authority to disqualify persons holding or seeking \emph{state} office. A five-Justice majority further held that responsibility for enforcing Section~3 against federal officeholders and candidates rests with Congress pursuant to its Section~5 enforcement power, requiring tailored federal legislation (\href{https://www.law.cornell.edu/supremecourt/text/23-719}{Cornell LII}).

\section{Reasoning}

The Per Curiam opinion grounded reversal in constitutional structure, text, Reconstruction history, and functional federalism, explicitly declining to reach whether Trump engaged in insurrection or whether the President is an ``officer of the United States.''

\textbf{1. The Fourteenth Amendment's Federalism Rebalancing.} Ratified in 1868, the Fourteenth Amendment ``expand[ed] federal power at the expense of state autonomy'' (\emph{Seminole Tribe of Fla. v. Florida}, 517 U.S. 44, 59 (1996); \emph{Fitzpatrick v. Bitzer}, 427 U.S. 445, 456 (1976)). While States retain traditional sovereign authority to prescribe qualifications for their \emph{own} state officers (\emph{Alden v. Maine}, 527 U.S. 706, 752 (1999); \emph{Taylor v. Beckham}, 178 U.S. 548, 570--571 (1900))---and enforced Section~3 against state officers after ratification (e.g., \emph{Worthy v. Barrett}, 63 N.C. 199 (1869))---powers over federal officers must be affirmatively ``delegated to, rather than reserved by, the States'' (\emph{U.S. Term Limits, Inc. v. Thornton}, 514 U.S. 779, 803--804 (1995)). Reading a state-restricting Amendment to silently empower States to disqualify federal candidates would invert its design (\href{https://www.law.cornell.edu/supremecourt/text/23-719}{Cornell LII}).

\textbf{2. Textual Limits and Interference with Congressional Amnesty.} Neither the Fourteenth Amendment nor the Elections and Electors Clauses (U.S.\ Const.\ art.\ I, \S 4, cl.\ 1; art.\ II, \S 1, cl.\ 2) delegates Section~3 enforcement over federal candidates to States. Moreover, Section~3's final sentence empowers Congress to remove any disability ``by a vote of two-thirds of each House'' at any time, including post-election (as Congress did in 1868 for Representative-elect Nelson Tift). Allowing States to bar federal candidates from the ballot before voting begins would burden Congress's amnesty power, violating the \emph{McCulloch v. Maryland}, 4 Wheat.\ 316, 436 (1819), principle that States may not impede federal constitutional operations.

\textbf{3. Historical Practice and Section~5 Exclusivity (Part II--A).} Invoking Chief Justice Chase's circuit opinion in \emph{Griffin's Case}, 11 F.\ Cas.\ 7, 26 (C.C.D.\ Va.\ 1869), and Senator Lyman Trumbull's 1869 floor statements, the five-Justice majority emphasized that because Section~3 imposes a ``preventive and severe penalty,'' formal ascertainment procedures are indispensable. Congress supplied such procedures via the Enforcement Act of 1870 (later repealed), the Confiscation Act of 1862 (now 18 U.S.C.\ \S 2383), and each House's Article~I, \S 5 power to judge its Members' qualifications. By contrast, no tradition existed of state enforcement against federal candidates (\emph{United States v. Texas}, 599 U.S. 670, 677 (2023)). Furthermore, because Section~5 is strictly remedial and subject to ``congruence and proportionality'' (\emph{City of Boerne v. Flores}, 521 U.S. 507, 520 (1997)), it is implausible that States---unbound by Section~5---would enjoy freer rein than Congress.

\textbf{4. Preventing a Chaotic Presidential ``Patchwork.''} Relying on \emph{Anderson v. Celebrezze}, 460 U.S. 780, 794--795 (1983), and \emph{U.S. Term Limits}, 514 U.S. at 822, the Court stressed that presidential elections implicate a ``uniquely important national interest.'' Divergent state standards of proof, evidentiary rules (such as admissibility of congressional reports), and ballot-challenge procedures (already producing exclusions in Colorado and Maine) would yield conflicting state-by-state outcomes, disrupt an evolving electoral map, and risk post-election chaos nullifying millions of votes (\href{https://www.law.cornell.edu/supremecourt/text/23-719}{Cornell LII}).

\section{Concurring and Dissenting Opinions}

\textbf{Justice Barrett (concurring in part and concurring in the judgment):} Joined Parts I and II--B of the Per Curiam opinion, agreeing that States cannot enforce Section~3 against presidential candidates. She declined to join the majority's Section~5 exclusivity analysis, noting that the suit arose under state law in state court and did not require deciding ``whether federal legislation is the exclusive vehicle through which Section~3 can be enforced.'' Urging institutional restraint during a volatile election season, she cautioned against amplifying judicial disagreement with stridency and underscored the Court's 9--0 unanimity on the outcome (\href{https://www.law.cornell.edu/supremecourt/text/23-719}{Cornell LII}).

\textbf{Justices Sotomayor, Kagan, and Jackson (jointly concurring in the judgment):} Agreed unanimously on federalism grounds that Colorado could not unilaterally exclude a presidential candidate and create a chaotic state-by-state patchwork. However, citing Chief Justice Roberts's restraint maxim in \emph{Dobbs v. Jackson Women's Health Organization}, 597 U.S. 215, 348 (2022), and Justice Breyer's dissent in \emph{Bush v. Gore}, 531 U.S. 98, 158 (2000), they sharply criticized the five-Justice majority for needlessly deciding how \emph{federal} enforcement must work. They argued that (1) Section~3's text is self-executing like other Reconstruction Amendments (\emph{City of Boerne}, 521 U.S. at 524) and presidential qualifications (Art.\ II, \S 1, cl.\ 5; Amdt.\ 22); (2) requiring a simple congressional majority to enact legislation undermines Section~3's two-thirds supermajority requirement for removing a disability; (3) \emph{Griffin's Case} is a nonprecedential circuit decision; and (4) the majority improperly foreclosed federal judicial enforcement or general statutory remedies to ``insulate this Court and petitioner from future controversy'' (\href{https://www.law.cornell.edu/supremecourt/text/23-719}{Cornell LII}).

\section{Key Quotations}

\begin{quote}
``Because the Constitution makes Congress, rather than the States, responsible for enforcing Section 3 against federal officeholders and candidates, we reverse.'' \\
--- Per Curiam, \href{https://www.law.cornell.edu/supremecourt/text/23-719}{\emph{Trump v. Anderson}, 601 U.S. 100 (2024)}
\end{quote}

\begin{quote}
``We conclude that States may disqualify persons holding or attempting to hold state office. But States have no power under the Constitution to enforce Section 3 with respect to federal offices, especially the Presidency.'' \\
--- Per Curiam, \href{https://www.law.cornell.edu/supremecourt/text/23-719}{\emph{Trump v. Anderson}, 601 U.S. 100 (2024)}
\end{quote}

\begin{quote}
``Nothing in the Constitution requires that we endure such chaos---arriving at any time or different times, up to and perhaps beyond the Inauguration.'' \\
--- Per Curiam, \href{https://www.law.cornell.edu/supremecourt/text/23-719}{\emph{Trump v. Anderson}, 601 U.S. 100 (2024)}
\end{quote}

\begin{quote}
``Particularly in this circumstance, writings on the Court should turn the national temperature down, not up. For present purposes, our differences are far less important than our unanimity: All nine Justices agree on the outcome of this case. That is the message Americans should take home.'' \\
--- Barrett, J., concurring in part and concurring in the judgment, \href{https://www.law.cornell.edu/supremecourt/text/23-719}{Cornell LII}
\end{quote}

\begin{quote}
``The Court today needed to resolve only a single question: whether an individual State may keep a Presidential candidate found to have engaged in insurrection off its ballot. The majority resolves much more than the case before us.'' \\
--- Sotomayor, Kagan, and Jackson, JJ., concurring in the judgment, \href{https://www.law.cornell.edu/supremecourt/text/23-719}{Cornell LII}
\end{quote}

\section{Doctrinal Significance and Modern Legacy (Through 2026)}

\emph{Trump v. Anderson} stands as a landmark structural federalism decision alongside \emph{U.S. Term Limits, Inc. v. Thornton} and \emph{Anderson v. Celebrezze}, cementing a sharp asymmetry in constitutional election administration: States enjoy broad retained authority over state-office qualifications (including Section~3 disqualifications, as in New Mexico's 2022 ouster of county commissioner Couy Griffin; \href{https://constitution.heritage.org/essays/amdt-14-s3/}{Heritage Guide}), yet possess zero unilateral power to enforce Section~3 against federal candidates (\href{https://constitution.congress.gov/browse/essay/intro.9-2-7/ALDE_00000048/}{Constitution Annotated}).

Immediately, the ruling overturned parallel ballot exclusions in Maine and Illinois, cleared Trump to appear on every 2024 primary and general-election ballot, and---paired with \emph{Trump v. United States}, 603 U.S. 593 (2024)---formed the legal backdrop to Trump's November 2024 election victory and January 2025 inauguration (\href{https://www.ussc.edu.au/constitutional-law-in-the-age-of-trump-2-0-the-supreme-court-s-2024-25-term-in-review}{USSC 2025 Review}). Through late 2026, Congress has enacted no new Section~5 civil enforcement statute, leaving 18 U.S.C.\ \S 2383 (criminal insurrection) and internal congressional seating powers under Article~I, \S 5 as the only active federal enforcement vehicles---rendering Section~3 largely dormant against federal candidates absent future congressional action.

In constitutional scholarship from 2024 through 2026, \emph{Trump v. Anderson} sparked an intense methodological debate. Prominent originalists William Baude and Michael Stokes Paulsen argued in \href{https://harvardlawreview.org/print/vol-138/sweeping-section-three-under-the-rug-a-comment-on-trump-v-anderson/}{138 Harv.\ L.\ Rev.\ 676 (2025)} that the Court ``swept Section Three under the rug'' by inverting Article~II's state-centric Electors Clause and abandoning originalism for pragmatic consequentialism. Similarly, Ilya Somin contended in the \href{https://www.cato.org/sites/cato.org/files/2024-09/cato-supreme-court-review-2023-2024-14.pdf}{Cato Supreme Court Review (2024)} that the Court ignored States' routine enforcement of Article~II age and citizenship qualifications (such as then-Judge Gorsuch's Tenth Circuit opinion in \emph{Hassan v. Colorado}). Meanwhile, state-vs.-federal ballot-access boundaries continue to generate certiorari petitions into the 2025--2026 Term testing how far \emph{Trump v. Anderson} restricts state election officials from adjudicating presidential-elector and candidate eligibility (\href{https://www.supremecourt.gov/DocketPDF/25/25-1244/407886/20260527090233541_20260527-085915-00001450-00003637.pdf}{No.\ 25--1244 Pet.\ for Cert.\ (2026)}).

\section{Pedagogical Critiques and Exam Hypotheticals}

\textbf{Key Critiques:} (1) \emph{Textual Inversion:} Requiring affirmative congressional legislation allows a bare majority (or one chamber's inaction) to nullify Section~3, contradicting its text requiring a two-thirds supermajority of both Houses to lift a disability. (2) \emph{Selective Originalism:} Critics note that the Court relied heavily on consequentialist ``patchwork'' fears and a single lower-court precedent (\emph{Griffin's Case}) rather than original public meaning. (3) \emph{Unresolved Merits:} The Court left unanswered whether the President is an ``officer of the United States,'' what conduct constitutes ``insurrection,'' and whether Congress could reject electoral votes on Section~3 grounds during a joint session without prior Section~5 legislation.

\textbf{``Opposite Side'' Contrast (State vs.\ Federal Qualification Enforcement):} Contrast \emph{Trump v. Anderson} with \emph{Chiafalo v. Washington}, 591 U.S. 578 (2020), and \emph{Hassan v. Colorado}, 495 F.\ App'x 947 (10th Cir.\ 2012) (\href{https://www.lawfaremedia.org/article/what-the-supreme-court-got-wrong-in-the-trump-section-3-case}{Lawfare}). In \emph{Chiafalo}, the Court upheld broad state power under Article~II's Electors Clause to penalize ``faithless'' presidential electors, and lower courts have routinely permitted States to exclude underage or non-natural-born presidential candidates from ballots. On an exam, students should distinguish objective, categorical Article~II qualifications (age, citizenship, residency) and state-controlled elector mechanics from Section~3's fact-intensive, Congress-waivable, Reconstruction-era disability---analyzing whether \emph{Trump v. Anderson}'s anti-patchwork rationale is strictly limited to Section~3 or threatens state enforcement of the Twenty-Second Amendment's two-term limit.

\textbf{Exam Hypothetical:} Suppose in 2028 a State Secretary of State excludes a former U.S.\ Senator from the U.S.\ Senate primary ballot after a state administrative hearing finds the candidate ``gave aid or comfort'' to a domestic rebellion, while simultaneously barring a former county sheriff from running for Governor on the exact same factual record. Under \emph{Trump v. Anderson}, the exclusion from the \emph{U.S.\ Senate} ballot is unconstitutional because States lack authority to enforce Section~3 against federal candidates; enforcement belongs exclusively to Congress (via Section~5 legislation or the Senate's Article~I, \S 5 judging power). By contrast, barring the former sheriff from the \emph{gubernatorial} ballot is constitutionally permissible under the Court's explicit carve-out preserving state power over state offices---subject, of course, to Fourteenth Amendment due-process constraints.

\end{document}
